\section{Procedencia y Adquisición del Movimiento}
\label{sec:procedencia}

El presente estudio investiga la viabilidad cinemática y la percepción ecológica de una posibilidad de acción (\emph{affordance}) en el tenis de mesa. Para ello, se emplea el registro biomecánico experimental de un gesto deportivo de ataque: el golpe \emph{topspin} (drive con efecto liftado) ejecutado con la extremidad superior izquierda.

\subsection{Fuente Experimental y Datos Xsens MVN}
El movimiento original fue adquirido mediante el sistema de captura inercial de cuerpo completo \Xsens{} (23 segmentos corporales, 17 sensores MTw Awinda) a una frecuencia de muestreo de $60.0\text{ Hz}$ ($16.67\text{ ms}$ entre fotogramas) en unidades métricas estandarizadas (metros para posición lineal y cuaterniones unitarios para orientación espacial). El ensayo analizado comprende un intervalo continuo de $6.300\text{ segundos}$ ($379\text{ fotogramas}$, numerados del $0$ al $378$), abarcando la fase de preparación estática inicial, el balanceo hacia atrás (\emph{backswing}), la aceleración explosiva, el impacto en máxima velocidad y la fase de desaceleración y seguimiento (\emph{follow-through}).

Los datos cinemáticos fuente provienen de los archivos tabulares estructurados:
\begin{itemize}[leftmargin=*]
  \item \texttt{Top\_spin\_SegmentOrientation.csv}: Contiene los cuaterniones normalizados $(q_0, q_1, q_2, q_3)$ para cada uno de los 23 segmentos anatómicos en cada instante de tiempo.
  \item \texttt{Top\_spin\_SegmentPosition.csv}: Registra las coordenadas tridimensionales $(x, y, z)$ del origen de cada segmento en el sistema de referencia inercial global de Xsens.
\end{itemize}

\subsection{Modelo Musculoesquelético en OpenSim}
Previamente al desarrollo en \Vizard{}, la cinemática inversa del gesto fue resuelta en \OpenSim{} \cite{delp2007opensim} utilizando el modelo musculoesquelético de cuerpo completo \texttt{Rajagopal2015.osim} \cite{rajagopal2016fullbody}, debidamente escalado antropométricamente a las dimensiones del atleta. De este procesamiento se obtuvieron las trayectorias articulares generalizadas (\texttt{Top\_spin-001\_Xsens\_jointangle\_q.mot}), así como las variables de referencia exportadas en los archivos \texttt{elbow\_flex\_l.sto} (ángulo de flexión del codo izquierdo) y \texttt{left\_hand\_height.sto} (coordenada de posición vertical $Y$ del centro de la mano izquierda). Estas trayectorias constituyen el patrón biomecánico fundamental para la validación numérica de fidelidad en la escena virtual.
